
\documentclass{article}
\usepackage{colortbl}
\usepackage{makecell}
\usepackage{multirow}
\usepackage{supertabular}

\begin{document}

\newcounter{utterance}

\centering \large Interaction Transcript for game `wordle', experiment `high\_frequency\_words\_no\_clue\_no\_critic', episode 4 with deepseek{-}v4{-}pro.
\vspace{24pt}

{ \footnotesize  \setcounter{utterance}{1}
\setlength{\tabcolsep}{0pt}
\begin{supertabular}{c@{$\;$}|p{.15\linewidth}@{}p{.15\linewidth}p{.15\linewidth}p{.15\linewidth}p{.15\linewidth}p{.15\linewidth}}
    \# & \multicolumn{2}{c}{Player} && \multicolumn{2}{c}{Game Master} \\
    \hline

    \theutterance \stepcounter{utterance}  
    & & & \multicolumn{4}{p{0.6\linewidth}}{
        \cellcolor[rgb]{0.9,0.9,0.9}{
            \makecell[{{p{\linewidth}}}]{
                \texttt{\tiny{[P1$\langle$GM]}}
                \texttt{Jesteś czarodziejem języka, który lubi zgadywać słowa, korzystając z podanych zasad.} \\
\\ 
\texttt{Witamy w Wordle! Masz sześć prób, aby odgadnąć słowo docelowe, poprawne polskie słowo składające się z pięciu małych liter. Użyj tagów "wyjaśnienie:" oraz "zgadnij:", aby podać zwięzłe wyjaśnienie dla każdej próby.} \\
\\ 
\texttt{Na przykład, jeśli Twoją próbą jest "dzień", Twoja odpowiedź powinna brzmieć} \\
\texttt{wyjaśnienie: to jest popularne pięcioliterowe polskie słowo i zaczynam zgadywanie od tego słowa.} \\
\texttt{zgadnij: dzień} \\
\\ 
\texttt{Po każdej próbie Twoja odpowiedź zostanie zweryfikowana i otrzymasz informację zwrotną wskazującą, które litery są poprawne (zielony), które litery są poprawne, ale na niewłaściwej pozycji (żółty), oraz które litery są niepoprawne (czerwony). Ta informacja zwrotna może być przydatna przy ustalaniu, które litery uwzględnić lub wykluczyć w następnej próbie.} \\
\\ 
\texttt{Na przykład informacja zwrotna dla "dzień" może wyglądać tak:} \\
\texttt{informacja zwrotna: d<żółty> z<żółty> i<zielony> e<żółty> ń<czerwony>} \\
\\ 
\texttt{Wyjaśnienie powinno zawierać szczegóły dotyczące tego, jak wykorzystano informacja zwrotna, aby dojść do nowej próby.} \\
\\ 
\texttt{Nie generuj żadnego innego tekstu.} \\
\\ 
\texttt{Zacznijmy od Twojej pierwszej próby.} \\
            }
        }
    }
     \\ \\

    \theutterance \stepcounter{utterance}  
    & \multicolumn{4}{p{0.6\linewidth}}{
        \cellcolor[rgb]{0.561,0.737,0.561}{
            \makecell[{{p{\linewidth}}}]{
                \texttt{\tiny{[P1$\rangle$GM]}}
                \texttt{wyjaśnienie: zaczynam od neutralnego słowa "kremy", które zawiera popularne samogłoski i spółgłoski, aby szybko sprawdzić obecność liter.} \\
\texttt{zgadnij: kremy} \\
            }
        }
    }
    & & \\ \\

    \theutterance \stepcounter{utterance}  
    & & & \multicolumn{2}{p{0.3\linewidth}}{
        \cellcolor[rgb]{0.9,0.9,0.9}{
            \makecell[{{p{\linewidth}}}]{
                \texttt{\tiny{[GM$|$GM]}}
                \texttt{attempts = 1} \\
\texttt{target\_word = układ} \\
\texttt{guess = kremy} \\
\texttt{guess\_feedback = k<yellow> r<red> e<red> m<red> y<red>} \\
            }
        }
    }
    & & \\ \\

    \theutterance \stepcounter{utterance}  
    & & & \multicolumn{4}{p{0.6\linewidth}}{
        \cellcolor[rgb]{0.9,0.9,0.9}{
            \makecell[{{p{\linewidth}}}]{
                \texttt{\tiny{[P1$\langle$GM]}}
                \texttt{informacja zwrotna: k<żółty> r<czerwony> e<czerwony> m<czerwony> y<czerwony>} \\
\\ 
\texttt{Podaj odpowiedź tylko w tym formacie.} \\
\texttt{wyjaśnienie: szczegóły} \\
\texttt{zgadnij: słowo} \\
            }
        }
    }
     \\ \\

    \theutterance \stepcounter{utterance}  
    & \multicolumn{4}{p{0.6\linewidth}}{
        \cellcolor[rgb]{0.561,0.737,0.561}{
            \makecell[{{p{\linewidth}}}]{
                \texttt{\tiny{[P1$\rangle$GM]}}
                \texttt{wyjaśnienie: Litera "k" jest w słowie, ale nie na pierwszej pozycji. "r", "e", "m", "y" są niepoprawne. Szukam słowa z "k" na innej pozycji, bez tych liter. Próbuję "lajki" – zawiera "k" na końcu i sprawdza nowe popularne litery.} \\
\texttt{zgadnij: lajki} \\
            }
        }
    }
    & & \\ \\

    \theutterance \stepcounter{utterance}  
    & & & \multicolumn{2}{p{0.3\linewidth}}{
        \cellcolor[rgb]{0.9,0.9,0.9}{
            \makecell[{{p{\linewidth}}}]{
                \texttt{\tiny{[GM$|$GM]}}
                \texttt{attempts = 2} \\
\texttt{target\_word = układ} \\
\texttt{guess = lajki} \\
\texttt{guess\_feedback = l<red> a<yellow> j<red> k<yellow> i<red>} \\
            }
        }
    }
    & & \\ \\

    \theutterance \stepcounter{utterance}  
    & & & \multicolumn{4}{p{0.6\linewidth}}{
        \cellcolor[rgb]{0.9,0.9,0.9}{
            \makecell[{{p{\linewidth}}}]{
                \texttt{\tiny{[P1$\langle$GM]}}
                \texttt{informacja zwrotna: l<czerwony> a<żółty> j<czerwony> k<żółty> i<czerwony>} \\
\\ 
\texttt{Podaj odpowiedź tylko w tym formacie.} \\
\texttt{wyjaśnienie: szczegóły} \\
\texttt{zgadnij: słowo} \\
            }
        }
    }
     \\ \\

    \theutterance \stepcounter{utterance}  
    & \multicolumn{4}{p{0.6\linewidth}}{
        \cellcolor[rgb]{0.561,0.737,0.561}{
            \makecell[{{p{\linewidth}}}]{
                \texttt{\tiny{[P1$\rangle$GM]}}
                \texttt{wyjaśnienie: Litery "a" i "k" są w słowie, ale nie na pozycjach 2 i 4. "l", "j", "i" są niepoprawne. Znamy już "a" i "k", brak "r", "e", "m", "y", "l", "j", "i". Szukam słowa z "a" i "k" na innych pozycjach, np. z "k" na początku lub końcu. Próbuję "kadet" – sprawdza "d", "e" (ponownie dla pewności), "t".} \\
\texttt{zgadnij: kadet} \\
            }
        }
    }
    & & \\ \\

    \theutterance \stepcounter{utterance}  
    & & & \multicolumn{2}{p{0.3\linewidth}}{
        \cellcolor[rgb]{0.9,0.9,0.9}{
            \makecell[{{p{\linewidth}}}]{
                \texttt{\tiny{[GM$|$GM]}}
                \texttt{attempts = 3} \\
\texttt{target\_word = układ} \\
\texttt{guess = kadet} \\
\texttt{guess\_feedback = k<yellow> a<yellow> d<yellow> e<red> t<red>} \\
            }
        }
    }
    & & \\ \\

    \theutterance \stepcounter{utterance}  
    & & & \multicolumn{4}{p{0.6\linewidth}}{
        \cellcolor[rgb]{0.9,0.9,0.9}{
            \makecell[{{p{\linewidth}}}]{
                \texttt{\tiny{[P1$\langle$GM]}}
                \texttt{informacja zwrotna: k<żółty> a<żółty> d<żółty> e<czerwony> t<czerwony>} \\
\\ 
\texttt{Podaj odpowiedź tylko w tym formacie.} \\
\texttt{wyjaśnienie: szczegóły} \\
\texttt{zgadnij: słowo} \\
            }
        }
    }
     \\ \\

    \theutterance \stepcounter{utterance}  
    & \multicolumn{4}{p{0.6\linewidth}}{
        \cellcolor[rgb]{0.561,0.737,0.561}{
            \makecell[{{p{\linewidth}}}]{
                \texttt{\tiny{[P1$\rangle$GM]}}
                \texttt{wyjaśnienie: Litery "k", "a", "d" są w słowie, ale nie na pozycjach 1, 2, 3. Z nieużytych jeszcze liter pasuje "o" i "n". Próbuję słowa "konda", gdzie "k" jest na 1, "a" na 5, "d" na 3 – to nie pasuje do wykluczonych pozycji? Sprawdzam: "k" nie na 1, "a" nie na 2, "d" nie na 3. W słowie "okład" – "k" na 2 (może być, bo żółte było na 1 i 4), "a" na 3 (może być, bo żółte na 2), "d" na 4 (może być, bo żółte na 3). Próbuję "okład".} \\
            }
        }
    }
    & & \\ \\

    \theutterance \stepcounter{utterance}  
    & & & \multicolumn{2}{p{0.3\linewidth}}{
        \cellcolor[rgb]{0.9,0.9,0.9}{
            \makecell[{{p{\linewidth}}}]{
                \texttt{\tiny{[GM$|$GM]}}
                \texttt{Error: The response should contain the 'zgadnij:' keyword exactly once.} \\
            }
        }
    }
    & & \\ \\

    \theutterance \stepcounter{utterance}  
    & & & \multicolumn{2}{p{0.3\linewidth}}{
        \cellcolor[rgb]{0.9,0.9,0.9}{
            \makecell[{{p{\linewidth}}}]{
                \texttt{\tiny{[GM$|$GM]}}
                \texttt{game\_result = ABORT} \\
            }
        }
    }
    & & \\ \\

\end{supertabular}
}

\end{document}
